\subsection{Filtered colimits, relative Theorem A, and strong saturation}
\label{sec:A5}

The two conclusions of XVI Corollary~7.7.2 have different requirements.
Finite forward excision already gives strong saturation.  Relative
Theorem~A requires a further infinitary condition: the counterexample below
satisfies even Loc~(9), but fails Loc~(8) and Loc~(3 bis).
Throughout this subsection ``weak localizer'' means precisely XVI
Loc~(1)--(3), including its cyclic saturation condition Loc~(1c).

\begin{claim}{A5.1}{Known}{Strong saturation needs neither Loc (8) nor Loc (9)}
\label{A51}
Every weak localizer satisfying Loc~($6'$a) is strongly saturated.
Thus the implication to Loc~(1 ter) in XVI Corollary~7.7.2 needs neither
Loc~(8) nor Loc~(9).
\end{claim}
\begin{proof}
We give the reduction to the exact saturation theorem used, including the
finite gluing step.  By \citet[Theorem~6.1.18]{Cisinski2006B},
$W_\infty\subseteq W$ under the weak localizer axioms alone.  Define
\[
 \mathcal V=\{f:X\to Y\text{ in }\mathrm{sSet}:\Delta/f\in W\}.
\]
It contains ordinary weak homotopy equivalences and satisfies two out of
three and the split-retract saturation condition: if $ri=1$ and $ir$ lies
in $\mathcal V$, then both $r,i$ lie in $\mathcal V$.  The latter is a
special case of Loc~(1c).

We check that monomorphisms in $\mathcal V$ are stable under pushouts.
For a monomorphism $X\hookrightarrow Y$ in $\mathcal V$ and any $X\to Z$,
write $T=Y\amalg_X Z$ and
$A=\Delta/X$, $B=\Delta/Y$, $C=\Delta/Z$.
Form the double collage $P$ with layers $A,B,C$.  In addition to the
arrows within the layers, it has arrows $a\to b$ indexed by
$\operatorname{Hom}_B(fa,b)$ and arrows $a\to c$ indexed by
$\operatorname{Hom}_C(ga,c)$, with the evident compositions, and no other
mixed arrows.  Its full subcategories $A\cup B$ and $A\cup C$ are sieves
covering $P$ and intersecting in $A$.  Retraction of either collage onto
its target layer is a left adjoint, hence a weak homotopy equivalence.
It follows that $A\to A\cup B$ lies in $W$.  Loc~($6'$a) gives
$A\cup C\to P\in W$, and hence $C\to P\in W$.

The canonical functor $P\to\Delta/T$ is an ordinary weak homotopy
equivalence.  Indeed, the two nerves of the collages are subcomplexes
covering $NP$ with intersection $NA$.  Their realization is the homotopy
pushout of $BB\leftarrow BA\to BC$.  The natural simplex-category
comparisons identify this with the homotopy pushout of
$|Y|\leftarrow|X|\to|Z|$, and the monomorphism $X\to Y$ makes the ordinary
simplicial pushout a homotopy pushout.  Consequently
$\Delta/Z\to\Delta/T\in W$.  This proves the required pushout closure
using finite forward excision only.

Now equip $\mathrm{sSet}$ with cofibrations the monomorphisms and weak
equivalences $\mathcal V$.  These satisfy exactly CF1--CF4 of
\citet[Definition~1.1.1]{RadulescuBanu2009}: the initial object and
composition axioms are immediate; CF2 is two out of three; CF3 is the
pushout assertion just proved.  For CF4 the ordinary mapping cylinder
factors every map as a monomorphism followed by an ordinary weak homotopy
equivalence, which belongs to $\mathcal V$.  All objects are cofibrant.
No arbitrary coproduct axiom CF5 or countable-composition axiom CF6 is
being asserted.  The split-retract condition is precisely the condition
WS3 in that source's Definition~7.2.6 (with its identities written with
the correctly typed source).  Its Theorem~7.2.7 therefore says that
$\mathcal V$ is saturated: a simplicial map is in $\mathcal V$ if and
only if it becomes invertible in
$\mathrm{sSet}[\mathcal V^{-1}]$.

For every category $D$, the natural last-vertex functor
$\Delta/ND\to D$ is an ordinary weak homotopy equivalence.  Its naturality
square and two out of three give
\[
                 f\in W\quad\Longleftrightarrow\quad Nf\in\mathcal V.
\]
Thus $N$ induces a functor
$\mathrm{Cat}[W^{-1}]\to\mathrm{sSet}[\mathcal V^{-1}]$.
If $f$ is inverted in the former localization, then $Nf$ is inverted in
the latter, so $Nf\in\mathcal V$ and $f\in W$.
This proves strong saturation before any use of a derivator or of relative
Theorem~A, and avoids the circularity identified by Grothendieck's
footnote~55 in XVI, ms.~118.
\end{proof}

\begin{claim}{A5.2}{Proved}{An accessible weak localizer with finite forward excision}
\label{A52}
For every functor $f:A\to B$ define its set of bad components by
\[
 E(f)=\{c\in\pi_0(BB): (BA)_c\longrightarrow (BB)_c
                   \text{ is not a weak homotopy equivalence}\},
\]
where the notation is used when $\pi_0(Bf)$ is bijective.  Let
$W_{\mathrm{fin}}$ consist of exactly those functors for which
\begin{enumerate}
\item $\Pi_1(Bf)$ is an equivalence of fundamental groupoids;
\item $H_*(Bf;\mathbb Z)$ is an isomorphism in every degree;
\item $E(f)$ is finite.
\end{enumerate}
This specifies membership for every functor of small categories.
The class satisfies XVI Loc~(1)--(3), Loc~($6'$a), Loc~(1 ter), and
Loc~(9).  Its aspherical categories are exactly the weakly contractible
categories.  In fact it is an $\aleph_1$-accessible full subcategory of
$\mathrm{Cat}^{[1]}$, and its inclusion preserves
$\aleph_1$-filtered colimits.
\end{claim}
\begin{proof}
First, identity maps and isomorphisms satisfy all three tests.  Fundamental
groupoid equivalences and integral homology isomorphisms satisfy two out
of three.  Once their component sets are identified, ordinary weak
equivalences satisfy two out of three component by component.  The bad
components of any one of $f,g,gf$ therefore lie in the union of the bad
components of the other two.  This proves two out of three for
$W_{\mathrm{fin}}$.

If $gf$ and $fg$ belong to $W_{\mathrm{fin}}$, their groupoid maps and
homology maps are isomorphisms or equivalences, so the maps of $f$ and
$g$ have the same property.  Away from the finitely many components
specified by the two composites, both composites are weak equivalences.
Ordinary weak homotopy equivalences satisfy two out of six, as is seen
on homotopy groups, or by the Whitehead theorem on their CW realizations.
Applied to the component maps of $f,g,f$ and of $g,f,g$, this makes
$f,g$ weak equivalences away from those finitely many components and
proves Loc~(1c).

A category with a final object is contractible, so satisfies Loc~(2).
More generally $A\to e$ lies in $W_{\mathrm{fin}}$ exactly when $BA$
is nonempty, simply connected, and integrally acyclic.  A simply connected
acyclic CW complex is weakly contractible by the Hurewicz theorem, applied
to its first possible nonzero homotopy group.  Conversely a weakly
contractible space plainly satisfies all three tests.  The condition on
bad components adds no restriction here.  If every comma category in
the absolute axiom Loc~(3) is $W_{\mathrm{fin}}$-aspherical, they are
therefore weakly contractible.  Quillen's Theorem~A, in its comma-category
dual form, makes the given functor an ordinary weak equivalence, hence a
member of $W_{\mathrm{fin}}$
\citep[Section~1, Theorem~A]{Quillen1973A5}.

We prove forward excision for \emph{every} homotopy pushout of spaces,
which includes all the MV squares required in Loc~($6'$a).  Consider
\[
\begin{tikzcd}
 X_0\ar[r,"a"]\ar[d,"i"']&X_2\ar[d,"j"]\\
 X_1\ar[r]&X_3 .
\end{tikzcd}
\]
Suppose $i$ satisfies the three defining tests.  Van Kampen's theorem
for fundamental groupoids shows that $\Pi_1(j)$ is an equivalence:
the groupoid homotopy pushout of an equivalence is an equivalence.
The integral homology Mayer--Vietoris sequence shows that $j$ is a homology
isomorphism.  Let $S\subseteq\pi_0(X_1)$ be the finite set of bad
components of $i$, and let $\bar S$ be its image in $\pi_0(X_3)$.
For a component $c\notin\bar S$, restrict the entire homotopy-pushout
diagram to the components mapping into $c$.  Its left vertical map is
now a coproduct of ordinary weak equivalences, so its right vertical map
is an ordinary weak equivalence.  Hence $E(j)\subseteq\bar S$ is finite.
For an MV cover of categories the two nerves are subcomplexes covering
the nerve of the union, so its classifying-space square is such a homotopy
pushout.  This proves Loc~($6'$a).  Strong saturation now follows from
A5.1, without being used in the preceding axiom checks.

We give the accessibility proof, since this is what prevents Loc~(9)
from substituting for Loc~(8).  Put $R=\operatorname{Ex}^{\infty}N$.
It takes countable categories to countable Kan complexes, preserves
monomorphisms and filtered colimits, and its homotopy groups calculate
those of the nerve.  To check the finitary assertions directly,
$(\operatorname{Ex}X)_n=\operatorname{Hom}(\operatorname{Sd}\Delta^n,X)$
and $\operatorname{Sd}\Delta^n$ is finite; $\operatorname{Ex}^{\infty}$
is the countable union of the iterates.  Fibrancy is
\citet[Proposition~2.1.39]{Cisinski2006B}.  Homotopy-group elements,
equalities and null homotopies in a Kan complex have finite simplicial
witnesses.  Cycles, boundaries and chain equalities also have finite
witnesses.  These observations give both of the following steps.

First let $f=\operatorname*{colim}_{i\in I}f_i$, where $I$ is
$\aleph_1$-filtered and each $f_i$ lies in $W_{\mathrm{fin}}$.
Fundamental-groupoid equivalences and integral homology isomorphisms
are stable under filtered colimits, by the finite-witness description
or, for homology, exactness of filtered colimits of abelian groups.
Suppose $f$ has infinitely many bad components.  Choose countably many
distinct ones.  In each choose a failure of injectivity or surjectivity
of some based homotopy-group map; the degree is at least two, since
$\Pi_1(f)$ is already an equivalence.  A failure of injectivity is
witnessed by a represented class together with a null homotopy of its
image; the class remains nonzero at every later stage because it is
nonzero in the colimit.  A failure of surjectivity is witnessed by a
represented target class; it cannot enter the image at any later stage
since it is not in the image in the colimit.  All representatives,
basepoints and null homotopies can be placed in one stage of the
$\aleph_1$-filtered diagram.  The components remain distinct there,
since their colimit components are distinct.  That stage would have
infinitely many bad components, a contradiction.  Thus $W_{\mathrm{fin}}$
is closed under $\aleph_1$-filtered colimits.

Second, every countable subfunctor of a fixed $f\in W_{\mathrm{fin}}$
is contained in a countable subfunctor $f':A'\to B'$ that belongs to
$W_{\mathrm{fin}}$.  Here subcategories need not be full.  Construct
countable subfunctors successively and take their countable union.
At every step add finite paths joining any two currently present objects
lying in the same ambient component, both in $A$ and in $B$, and add a
source object and an image-to-object path for each currently present
target object.  This makes the eventual component sets inject into the
ambient component sets and makes $\pi_0(f')$ bijective.  Add witnesses
for surjectivity on $\pi_1$ and homology, and witnesses killing kernels
whenever a kernel witness is already present.  In an ambient component
where $f$ is an ordinary weak equivalence, do the same for every
$\pi_n$, $n\geq2$.  These witnesses exist in the ambient categories
because $f$ satisfies the relevant tests there.  They are countable in
number: at a countable stage there are countably many finite simplicial
or integral-chain tests.  Each witness uses finitely many simplices of
$R(A)$ or $R(B)$ and therefore countably, in fact finitely up to generated
subcategory, many ambient arrows.  Close under composition and under
the image functor, still adding only countably many arrows.

Every test occurring in the union occurs at a finite stage and its
required witnesses occur later.  The resulting $f'$ is therefore a
fundamental-groupoid equivalence and a homology isomorphism.  On each
ambient good component it is a weak equivalence.  It has at most one
component within each ambient component, by the path additions, and so
has at most the finitely many ambient bad components as exceptions.
Thus $f'\in W_{\mathrm{fin}}$.

The poset of countable subfunctors of $f$ that lie in
$W_{\mathrm{fin}}$ is consequently $\aleph_1$-filtered: any countable
family has countable union and the preceding construction extends that
union to another member.  Its colimit is $f$.  A category with countably
many arrows, and an arrow between two such categories, is
$\aleph_1$-presentable: it has a countable generators-and-relations
presentation, whose data and equations lift through an
$\aleph_1$-filtered diagram.  Since the full inclusion of
$W_{\mathrm{fin}}$ preserves these colimits, its countable objects are
$\aleph_1$-presentable in that full subcategory as well.  There is a set
of their isomorphism classes.  They generate it under
$\aleph_1$-filtered colimits, which proves the asserted accessibility
and Loc~(9).
\end{proof}

\begin{claim}{A5.3}{Refuted}{Relative Theorem A without filtered-colimit closure}
\label{A53}
The implication
\[
 \mathrm{Loc}(1)\text{--}(3),\ \mathrm{Loc}(6'\mathrm a),\
 \mathrm{Loc}(9)\quad\Longrightarrow\quad\mathrm{Loc}(3\mathrm{bis})
\]
is false.  For $W_{\mathrm{fin}}$ its failure occurs over the discrete
category $\mathbb N$, with every comma functor equal to one explicit
functor $j:Q\to P$ between finite categories.  Here $Q$ has two objects
and four arrows; $P$ has 58 objects and 518 arrows.
\end{claim}
\begin{proof}
We give the finite categories completely.  Let $C_n$ be the polygonal
circle with vertices $0,\ldots,n-1$ and edges
$\{i,i+1\pmod n\}$.  Let $S_n=C_n*\{N,S\}$ be its suspension,
a triangulated 2-sphere.  Write $L$ and $K$ for the posets of nonempty
faces of $S_6$ and $S_3$, ordered by inclusion.  Define simplicial maps
$u,v:S_6\to S_3$ fixing the two poles and sending the equatorial
vertices, respectively, to
\[
 u=(0,0,1,1,2,2),\qquad v=(0,1,2,0,1,2).
\]
Their induced face-poset maps, also denoted $u,v:L\to K$, send a face
to its nonempty image.  The maps on the equatorial circles have degrees
one and two, so their suspended maps have those degrees on $H_2$.

The objects of $P$ are the disjoint union of the objects of $L$ and $K$.
Within each layer the arrows are its poset arrows.  For $a\in L$,
$b\in K$, the mixed arrows $a\to b$ are
\[
 \{(0,a,b):u(a)\subseteq b\}\ \amalg\
 \{(1,a,b):v(a)\subseteq b\}.
\]
There are no arrows from $K$ to $L$.  Composition with an arrow in
either layer retains the mixed arrow's label 0 or 1.  These rules
specify all arrows, identities and compositions.  Let $Q$ be the full
subcategory on the two north-pole faces, one in each layer.  It has
two parallel nonidentity arrows from the first pole to the second.
Let $j:Q\hookrightarrow P$ be the inclusion.

The realization $BP$ is the homotopy coequalizer of
$Bu,Bv:BL\rightrightarrows BK$.  One can verify this without a
general homotopy-colimit theorem: the nerves of the two single-label
collages cover $NP$, with intersection $NL\amalg NK$; each collage
is a mapping-cylinder model for its indicated functor, relative to its
source and target.  Their union is the double mapping cylinder.  Since
$BL$ and $BK$ are 2-spheres and $Bu$ has degree one, this is homotopy
equivalent to the mapping torus of a based degree-two map of $S^2$.
The pole subcategory is its section circle.  Van Kampen therefore gives
\[
                \pi_1(BQ)\xrightarrow{\cong}\pi_1(BP)\cong\mathbb Z.
\]
Equivalently, collapsing a maximal tree in the double cylinder leaves
its one circle generator; both sphere layers are simply connected and
add no fundamental-group relation.

Put $R=\mathbb Z[t,t^{-1}]$, where $t$ is this circle generator.
The relative universal-cover chain complex, up to chain homotopy, is
\[
 C_*(\widetilde{BP},\widetilde{BQ})\simeq
       \bigl[R\xrightarrow{\,2t-1\,}R\bigr],
              \qquad\text{in degrees }3,2.
\]
Indeed, cut the mapping torus along its sphere fibre and lift it to
the line: each adjacent pair of sphere copies is attached by the
degree-one and degree-two ends, giving $t v_*-u_*$ on the reduced
sphere chains.  This argument also shows that a finite CW model of the
pair is $(S^1\vee S^2)\cup_{2t-1}e^3$ relative to $S^1$.
After augmentation $t=1$ the differential is $1$.  Thus $Bj$ is an
integral homology equivalence.  But
\[
 H_2(\widetilde{BP},\widetilde{BQ})
    =R/(2t-1)\cong\mathbb Z[1/2]\ne0,
\]
where $t$ acts by $1/2$.  The map $Bj$ is not a weak equivalence.
Both categories are connected, so $j\in W_{\mathrm{fin}}$ with exactly
one bad component.

The attached program
\texttt{code/A5\_finite\_component\_counterexample.py} constructs all
58 objects and 518 arrows and checks all 1,920 composable pairs and
4,924 composable triples.  It computes the two degrees on explicit
integral fundamental cycles and the homology of the double-cylinder
chain complex, whose chain ranks are $(5,17,24,12)$ and differential
ranks $(4,12,12)$.  Smith normal form gives $H_0=H_1=\mathbb Z$ and
zero higher homology, with no torsion.  It also checks exactly that
$2t-1$ augments to 1 and vanishes under the nonzero homomorphism
$R\to\mathbb Z[1/2]$, $t\mapsto1/2$.  The program is a finite
certificate for the stated categories and attaching calculation;
the infinite argument below is mathematical, not a claimed exhaustive
machine enumeration.

Let $D$ be the discrete category on $\mathbb N$ and form
\[
 f=\coprod_{n\in\mathbb N}j:
       \coprod_{n\in\mathbb N}Q\longrightarrow
       \coprod_{n\in\mathbb N}P
\]
over $D$, each summand being sent to its index.  For every $n\in D$,
the comma functor $f/n$ is exactly $j$, hence belongs to
$W_{\mathrm{fin}}$.  The map $f$ is a groupoid equivalence and an
integral homology equivalence, but it has infinitely many bad components.
Consequently $f\notin W_{\mathrm{fin}}$, disproving relative
Theorem~A.  Each finite coproduct $\coprod_{n<m}j$ lies in
$W_{\mathrm{fin}}$, and their sequential filtered colimit is $f$.
Thus Loc~(8) and even Loc~(8 bis) fail.  This completes the verification
of every asserted axiom and every failure on all small categories.
\end{proof}

\begin{claim}{A5.4}{Known}{The corrected implication and the manuscript footnotes}
\label{A54}
Under XVI Loc~(1)--(3) and Loc~($6'$a), Loc~(8) suffices for relative
Theorem~A, without Loc~(9).  Loc~(9) alone does not suffice.
For strong saturation neither is needed.  Thus Lead~1 is true for the
relative conclusion, and false if read as asserting a need for Loc~(8)
for the saturation conclusion.
\end{claim}
\begin{proof}[Audit of the positive implication]
We read DER\_FOLLOWUP G.9 and checked its exact point of dependence on
filtered colimits.  The finite double-collage argument in A5.1 makes
monomorphisms in $\mathcal V$ stable under pushouts.  The functor
$X\mapsto\Delta/X$ preserves filtered colimits, since its objects,
arrows and equations are finite data.  Loc~(8) therefore gives the
limit-ordinal steps in transfinite composites of monomorphisms in
$\mathcal V$.  Together with two out of three and the inclusion of
maps right orthogonal to all monomorphisms, these are precisely
Cisinski's $\Delta$-localizer axioms.  His Proposition~3.4.25 makes
every such localizer regular; its representables are contractible,
so it is a test localizer.  His Theorem~4.2.15 associates a fundamental
localizer $W'$ on categories.  The natural weak equivalences
$\Delta/NA\to A$ identify
\[
       W=N^{-1}\mathcal V=W'
\]
by two out of three.  Hence $W$ satisfies relative Theorem~A.
This argument uses no accessibility theorem.  The exact thesis/book
numbering comparison is B.8; the book's weak-minimality theorem
6.1.18 cannot be silently replaced by the thesis's stronger-premise
minimality statement.

A5.2--A5.3 show that keeping accessibility while deleting
Loc~(8) does not recover the relative axiom.  A5.1 proves the separate
saturation assertion.  Both derivations avoid identifying a weak
localizer with the localizer of its purported derivator before the
relative axiom is established.  Thus they avoid the circle named in
Grothendieck's footnote~55.  His footnote~56, ``Loc (3 bis) pas prouv\'e,
probablement faux'', is incorrect for the stated hypotheses
Loc~($6'$a,8,9), but is vindicated after Loc~(8) is removed, even with
Loc~(9) retained.  These are Grothendieck's own footnotes, not notes
supplied by the edition.
\end{proof}
